\section{Discussion}\label{sec:discussion}

\paragraph{Conventional and proof-based ranks.}
The proof-based sizes are about a quarter larger than the conventional ones. The difference is the slack of the projection, which multiplies all norms of the recursive levels, the factor $b+1$ for the decomposed opening, and above all the norm for the inner commitments of the terminal level: $8T\bar\gamma$ in the proof, $2\bar\gamma$ in the convention that reproduces the published sizes. The following observation shows that the convention cannot be justified by a proof that goes through the uniqueness of weak openings.

\begin{proposition}\label{prop:weak}
Let $A\in\Rq^{\kappa\times m}$ and $\nu>0$, and let $\mathbf x\in R^m$ be non-zero modulo $q$ with $A\mathbf x=0\bmod q$ and $\norm{\mathbf x}\le\frac{25}{4}\bigl(\nu-2\sqrt{md}\bigr)$. Then one can compute from $\mathbf x$ a vector $\vt\in\Rq^\kappa$ and two weak openings $(\bar\vs,0,\delta)$ and $(\bar\vs',0,\delta')$ of $\vt$ of norm $\nu$, for $\varrho=0$, with $\bar\vs\ne\bar\vs'$.
\end{proposition}
\begin{proof}
Let $\delta=3(1-X)$ and $\delta'=4(1-X)$. Both are differences of two elements of $\Ch$: take a challenge with coefficients $2,-1$, respectively $2,-2$, at $X^0$ and $X^1$ and exchange the two coefficients; \cref{app:pairs} lists two such pairs in which all four polynomials have operator norm below $15$. Let $\mathbf w=-\lfloor7\mathbf x/25\rceil$, rounding coefficient-wise, $\mathbf u=\mathbf x+3\mathbf w$ and $\mathbf u'=\mathbf x+4\mathbf w$. Then $4\mathbf u-3\mathbf u'=\mathbf x$, and the coefficients of $\mathbf u-\frac4{25}\mathbf x$ and of $\mathbf u'+\frac3{25}\mathbf x$ are bounded by $3/2$ and $2$ in absolute value, so that $\norm{\mathbf u}\le\frac4{25}\norm{\mathbf x}+\frac32\sqrt{md}\le\nu$ and $\norm{\mathbf u'}\le\frac3{25}\norm{\mathbf x}+2\sqrt{md}\le\nu$. Put $\bar\vs=\delta^{-1}\mathbf u$ and $\bar\vs'=\delta'^{-1}\mathbf u'$ in $\Rq^m$ and $\vt=A\bar\vs$. Then $\bar\vs-\bar\vs'=(12(1-X))^{-1}(4\mathbf u-3\mathbf u')=(12(1-X))^{-1}\mathbf x$, which is non-zero and in the kernel of $A$. Hence $A\bar\vs'=\vt$, and $\delta\bar\vs=\mathbf u$ and $\delta'\bar\vs'=\mathbf u'$ have norm at most $\nu$.
\end{proof}

So weak openings of norm $\nu$ are unique only if Module-SIS is hard for norms up to about $6\nu$; \cref{lem:binding}(iii) shows that hardness for $4T\nu=60\nu$ suffices. In the terminal level a weak opening has norm $\nu=2\bar\gamma$. With ranks chosen for the norm $2\bar\gamma$, the assumption does not exclude kernel vectors of norm $12\bar\gamma$, from which an adversary obtains inner commitments with several weak openings, whereas the proofs of~\cite{BS23}, of~\cite{NS24} and of \cref{thm:level} all use that the prover is bound to vectors of this kind. We know of no attack on the protocols with conventional ranks, and of no proof. Closing the gap from either side, by a proof with a smaller constant than $8T$ or by an attack, would change the sizes of all protocols of this family by several kilobytes.

\paragraph{Verifier and commitment.}
The tail inherits from LaBRADOR a verifier that is linear in the statement, and \cref{sec:pcs} inherits from Greyhound a statement of size $\sqrt N$ and a commitment of more than a kilobyte. Akita verifies in milliseconds, $22$ to $109$ times faster than Greyhound, and its commitments have $128$ bytes~\cite{DBK26}. Our result concerns the proof size only. Whether the tail can follow a first fold with a succinct verifier without losing that property depends on the structure of the constraint vectors: the verifier of a level has to evaluate the multilinear extensions $\tilde\Phi_i(\rho)$, which is cheap if the constraints of the statement are structured, as the rows of~\eqref{eq:gh-system} are, but the commitment keys $A$, $B$, $D$ of the previous level are not.

\paragraph{The non-interactive protocol.}
\Cref{thm:level,thm:pcs} are about interactive protocols, and \cref{sec:fs} explains why the known results on the Fiat--Shamir transformation of multi-round protocols do not apply to protocols of this family. A proof of knowledge soundness in the random oracle model, classical or quantum, for LaBRADOR-style arguments is in our view the most important missing piece of the security analysis; the sum-check adds $\mu$ rounds per level to the rounds of LaBRADOR but does not change the nature of the problem.

\paragraph{Zero knowledge.}
The protocol is not zero-knowledge: $\mathbf p$, the lifted values, the sum-check messages and the terminal opening depend on the witness. The techniques that make LaBRADOR zero-knowledge~\cite{BS23,BBLMOSS26} mask the witness and the projection; the sum-check messages would need masking in addition, for example by a committed random polynomial. We have not worked this out.

\paragraph{Completeness and real witnesses.}
The guaranteed parameters make completeness a theorem, at a cost of about a sixth of the proof. The optimised parameters rely, like those of LaBRADOR and Greyhound, on expected norms, and the end-to-end runs of \cref{sec:impl} show where this breaks for polynomial commitments: not in the tail, whose witnesses are close to their expected norms for every polynomial we tried, but in the first fold. If many blocks of the polynomial have the same digits, the amortised opening $\sum_ic_i\vs_i$ contains $(\sum_ic_i)\vs_1$, and the share of the norm that depends on it grows with $N$; for a constant polynomial with $2^{28}$ coefficients the honest prover failed. Commitments to structured data, for instance vectors with many repeated entries, are common in applications, so a deployment should use bounds that hold for every polynomial at least in the first fold. In our runs the guaranteed proofs were only \nEeGcBelowLo\% to \nEeGcBelowHi\% below their bounds, so the cost of guaranteed completeness is close to the difference between the settings in the model. Tighter bounds than those of \cref{thm:complete}, for example from the variance of $\norm{\vz}^2$, would narrow this gap; the runs suggest that such bounds must allow for heavy tails when blocks repeat.

\paragraph{Other moduli and rings.}
The construction needs a field between $\Z_q$ and $\Rq$ with about $2^{128}$ elements, which requires that the irreducible factors of $X^d+1$ modulo $q$ have degree at least $128/\log_2q$; for $q\equiv5\pmod8$ every power of two $f\le d/2$ is available. For moduli with fast number-theoretic transforms one would have to work with a field that is not contained in the ring, and the lifting of \cref{sec:field} does not apply as it stands.
